% Part: turing-machines
% Chapter: machines-computations
% Section: representing-tms

\documentclass[../../../include/open-logic-section]{subfiles}

\begin{document}

\olfileid{tur}{mac}{rep}
\olsection{تمثيل آلات تورنغ}

\begin{explain}
يمكن تمثيل آلات تورنغ بصريًا بواسطة \emph{مخططات الحالات}. تتألف هذه
المخططات من عقد حالات تصل بينها أسهم. وتمثل كل عقدة، كما هو متوقع،
حالة من حالات الآلة. ويمثل كل سهم تعليمة يمكن تنفيذها انطلاقًا من تلك
الحالة، وتكتب تفاصيل التعليمة فوق السهم المناسب أو تحته. انظر إلى
الآلة الآتية، التي ليس لها إلا حالتان داخليتان، $q_0$ و$q_1$، وتعليمة
واحدة:
\[
\begin{tikzpicture}[->,>=stealth',shorten >=1pt,auto,node distance=2.8cm,
                    semithick]
  \tikzstyle{every state}=[fill=none,draw=black,text=black]

  \node[initial,state]   (A)              {$q_0$};
  \node[state]   (B) [right of=A] {$q_1$};

  \path (A) edge  node {\TMtrans{\TMblank}{\TMstroke}{\TMright}} (B);
\end{tikzpicture}
\]
تذكر أن لآلة تورنغ رأسًا للقراءة والكتابة وشريطًا كتب عليه الدخل.
وتقرأ التعليمة هكذا: \emph{إذا كنت تقرأ $\TMblank$ في الحالة $q_0$،
فاكتب $\TMstroke$ وتحرك يمينًا وانتقل إلى الحالة $q_1$}. وهذا يكافئ
أن ترسل دالة الانتقال $\tuple{q_0,\TMblank}$ إلى
$\tuple{q_1,\TMstroke,\TMright}$.
\end{explain}

\begin{ex}
\emph{آلة الزوجية}: تتوقف آلة تورنغ الآتية إذا وفقط إذا كان على
الشريط عدد زوجي من رموز $\TMstroke$ (بافتراض أن كل رموز $\TMstroke$
تسبق أول $\TMblank$ على الشريط).
\[
\begin{tikzpicture}[->,>=stealth',shorten >=1pt,auto,node distance=2.8cm,
                    semithick]
  \tikzstyle{every state}=[fill=none,draw=black,text=black]

  \node[initial,state]         (A)              {$q_0$};
  \node[state]         (B) [right of=A] {$q_1$};

  \path (A) edge [bend left] node {\TMtrans{\TMstroke}{\TMstroke}{\TMright}} (B)
        (B) edge [loop above] node {\TMtrans{\TMblank}{\TMblank}{\TMright}} (B)
            edge [bend left] node {\TMtrans{\TMstroke}{\TMstroke}{\TMright}} (A);
\end{tikzpicture}
\]

يقابل مخطط الحالات دالة الانتقال الآتية:
\begin{align*}
\delta(q_0, \TMstroke) & = \tuple{q_1, \TMstroke, \TMright},\\
\delta(q_1, \TMstroke) & = \tuple{q_0, \TMstroke, \TMright},\\
\delta(q_1, \TMblank)  & = \tuple{q_1, \TMblank, \TMright}
\end{align*}
\end{ex}

\begin{explain}
لا تتوقف الآلة السابقة إلا عندما يكون الدخل عددًا زوجيًا من العلامات.
وإلا واصلت الآلة، من الناحية النظرية، عملها إلى غير نهاية. ويمكن، لأي
آلة ودخل، تتبع \emph{تهيئات} الآلة لتحديد الخرج. وسنقدم تعريفًا صوريًا
للتهيئات لاحقًا. أما الآن فيمكن أن ننظر إليها حدسيًا بوصفها سلسلة من
الرسوم تبين حالة الآلة في أي لحظة أثناء التشغيل. وتبين التهيئات محتوى
الشريط، وحالة الآلة، وموضع رأس القراءة والكتابة.

لنتتبع تهيئات آلة الزوجية حين تبدأ بدخل مؤلف من أربع علامات
$\TMstroke$. نتوقع في هذه الحال أن تتوقف الآلة. ثم سنشغلها على دخل
مؤلف من ثلاث علامات $\TMstroke$، فستعمل إلى الأبد.

تبدأ الآلة في الحالة~$q_0$، وهي تمسح علامة~$\TMstroke$ الواقعة في
أقصى اليسار. ويمكننا تمثيل التهيئة الابتدائية للآلة كما يأتي:
\[
\TMendtape \TMstroke_0 \TMstroke \TMstroke \TMstroke \TMblank \ldots
\]
هذه التهيئة واضحة. فالآلة تبدأ في الحالة~$q_0$، وتمسح
علامة~$\TMstroke$ الواقعة في أقصى اليسار. ويمثل ذلك بوضع اسم الحالة
دليلًا سفليًا على أول علامة~$\TMstroke$. والتعليمة المنطبقة عند هذه
النقطة هي
$\delta(q_0,\TMstroke)=\tuple{q_1,\TMstroke,\TMright}$، ولذلك تتحرك
الآلة يمينًا على الشريط وتنتقل إلى الحالة~$q_1$.
\[
\TMendtape \TMstroke \TMstroke_1 \TMstroke \TMstroke \TMblank \ldots
\]
ولما كانت الآلة الآن في الحالة~$q_1$ وتمسح~$\TMstroke$، فعلينا
«اتباع» التعليمة $\delta(q_1,\TMstroke)=\tuple{q_0,
  \TMstroke,\TMright}$. وينتج من ذلك التهيئة
\[
\TMendtape \TMstroke \TMstroke \TMstroke_0 \TMstroke \TMblank \ldots
\]
ومع استمرار الآلة تطبق القواعد بالترتيب نفسه مرة أخرى، فتنتج
التهيئتان الآتيتان:
\[
\TMendtape \TMstroke \TMstroke \TMstroke \TMstroke_1 \TMblank \ldots
\]
\[
\TMendtape \TMstroke \TMstroke \TMstroke \TMstroke \TMblank_0 \ldots
\]
الآلة الآن في الحالة~$q_0$ وتمسح~$\TMblank$. ومن مخطط الانتقال نرى
بسهولة أنه لا توجد تعليمة تنفذ، ولذلك توقفت الآلة. وهذا يعني أن الدخل
قد قُبل.

ولنفترض بعد ذلك أننا بدأنا الآلة بدخل مؤلف من ثلاث علامات
$\TMstroke$. تتشابه التهيئات الأولى لأن التعليمات نفسها تنفذ، مع فرق
صغير في دخل الشريط:
\[
\TMendtape \TMstroke_0 \TMstroke \TMstroke \TMblank \ldots
\]
\[
\TMendtape \TMstroke \TMstroke_1 \TMstroke \TMblank \ldots
\]
\[
\TMendtape \TMstroke \TMstroke \TMstroke_0 \TMblank \ldots
\]
\[
\TMendtape \TMstroke \TMstroke \TMstroke \TMblank_1 \ldots
\]
لقد اجتازت الآلة الآن جميع علامات $\TMstroke$، وهي تقرأ~$\TMblank$
في الحالة~$q_1$. وكما يبين المخطط، توجد تعليمة من الصورة
$\delta(q_1,\TMblank)=\tuple{q_1,\TMblank,\TMright}$. ولما كان الشريط
مملوءًا إلى غير نهاية برموز $\TMblank$ نحو اليمين، فستواصل الآلة
تنفيذ هذه التعليمة \emph{إلى الأبد}، باقية في الحالة~$q_1$ ومتحركة
مسافة متزايدة نحو اليمين. ولن تتوقف الآلة قط، فلا تقبل الدخل.
\end{explain}

\begin{explain}
من المهم ملاحظة أن الآلات لا تتوقف كلها. فإذا كان التوقف يعني نفاد
التعليمات التي تنفذها الآلة، أمكننا إنشاء آلة لا تتوقف قط بمجرد ضمان
وجود سهم خارج لكل رمز في كل حالة. ويمكن تعديل آلة الزوجية كي تعمل
إلى غير نهاية بإضافة تعليمة لمسح $\TMblank$ في~$q_0$.
\end{explain}

\begin{ex}
\[
\begin{tikzpicture}[->,>=stealth',shorten >=1pt,auto,node distance=2.8cm,
                    semithick]
  \tikzstyle{every state}=[fill=none,draw=black,text=black]

  \node[initial,state] (A)              {$q_0$};
  \node[state]         (B) [right of=A] {$q_1$};

  \path
  (A) edge [bend left] node {\TMtrans{\TMstroke}{\TMstroke}{\TMright}} (B)
      edge [loop above] node {\TMtrans{\TMblank}{\TMblank}{\TMright}} (A)
  (B) edge [loop above] node {\TMtrans{\TMblank}{\TMblank}{\TMright}} (B)
      edge [bend left] node {\TMtrans{\TMstroke}{\TMstroke}{\TMright}} (A);
\end{tikzpicture}
\]
\end{ex}

\begin{explain}
جداول الآلة طريقة أخرى لتمثيل آلات تورنغ. تعرض هذه الجداول أبجدية
الشريط على المحور $x$، ومجموعة حالات الآلة على المحور $y$. ويكتب داخل
الجدول، عند تقاطع كل حالة ورمز، باقي التعليمة: الحالة الجديدة والرمز
الجديد واتجاه الحركة. وتسهل جداول الآلة تحديد الحالة والرمز اللذين
تتوقف عندهما الآلة. فكل خانة فارغة في الجدول موضع محتمل لتوقف الآلة.
وفي حين لا تتضح مواضع توقف الآلة دائمًا من مخططات الحالات ومجموعات
التعليمات مباشرة، يسهل التعرف على جميع مواضع التوقف بالعثور على
الخانات الفارغة في جدول الآلة.
\end{explain}

\begin{ex}
جدول آلة الزوجية هو:
\[
\centering
\begin{tabular}{lllll}
\cline{2-4}
\multicolumn{1}{l|}{}      & \multicolumn{1}{c|}{$\TMblank$}
& \multicolumn{1}{c|}{$\TMstroke$}     & \multicolumn{1}{c|}{$\TMendtape$}   &  \\ \cline{1-4}
\multicolumn{1}{|l|}{$q_0$} & \multicolumn{1}{l|}{}
& \multicolumn{1}{l|}{\TMtrans{\TMstroke}{q_1}{\TMright}} &
\multicolumn{1}{l|}{\phantom{\TMtrans{\TMstroke}{q_1}{\TMright}}} &  \\ \cline{1-4}
\multicolumn{1}{|l|}{$q_1$} & \multicolumn{1}{l|}{\TMtrans{\TMblank}{q_1}{\TMright}}
& \multicolumn{1}{l|}{\TMtrans{\TMstroke}{q_0}{\TMright}} &
\multicolumn{1}{l|}{\phantom{\TMtrans{\TMstroke}{q_1}{\TMright}}} &  \\ \cline{1-4}
\end{tabular}
\]
وكما نرى، تتوقف الآلة حين تمسح~$\TMblank$ في الحالة~$q_0$.
\end{ex}

\begin{explain}
لم ننظر حتى الآن إلا في آلات تقرأ الدخل وتقبله. لكن آلات تورنغ قادرة
على القراءة والكتابة معًا. ومن أمثلة هذه الآلات (وهي كثيرة جدًا)
\emph{آلة المضاعفة}. فإذا بدأت آلة المضاعفة بكتلة من $n$ من علامات
$\TMstroke$ على الشريط، أخرجت كتلة من $2n$ من علامات $\TMstroke$.
\end{explain}

\begin{ex}
  \ollabel{ex:doubler} قبل بناء آلة مضاعفة، من المهم وضع
  \emph{استراتيجية} لحل المسألة. ولما كانت الآلة (بالصيغة التي
  قدمناها) لا تستطيع تذكر عدد علامات $\TMstroke$ التي قرأتها، فعلينا
  ابتكار طريقة لتتبع كل علامات $\TMstroke$ على الشريط. ومن هذه الطرائق
  فصل الخرج عن الدخل بعلامة~$\TMblank$. وعندئذ تستطيع الآلة محو أول
  $\TMstroke$ من الدخل، واجتياز بقية الدخل، وترك $\TMblank$، وكتابة
  علامتي $\TMstroke$ جديدتين. ثم تعود الآلة وتجد ثاني $\TMstroke$ في
  الدخل وتضاعفها أيضًا. فهي تكتب، لكل $\TMstroke$ واحدة من الدخل،
  علامتي $\TMstroke$ من الخرج. وبمحو الدخل أثناء سير الآلة نضمن ألا
  تفوت أي علامة $\TMstroke$ وألا تضاعف أي منها مرتين. وحين يمحى الدخل
  كله تبقى على الشريط $2n$ من علامات $\TMstroke$. ويصور
  \olref{fig:doubler} مخطط حالات آلة تورنغ الناتجة.
  \begin{figure}
\[
\begin{tikzpicture}[->,>=stealth',shorten >=1pt,auto,node distance=2.8cm,
                    semithick]
  \tikzstyle{every state}=[fill=none,draw=black,text=black]

  \node[initial,state] (1)              {$q_0$};
  \node[state]         (2) [right of=1] {$q_1$};
  \node[state]         (3) [right of=2] {$q_2$};
  \node[state]         (4) [below of=3] {$q_3$};
  \node[state]         (5) [left of=4]  {$q_4$};
  \node[state]         (6) [left of=5]  {$q_5$};

  \path (1) edge node {\TMtrans{\TMstroke}{\TMblank}{\TMright}} (2)
    (2) edge [loop above] node {\TMtrans{\TMstroke}{\TMstroke}{\TMright}} (2)
      edge node {\TMtrans{\TMblank}{\TMblank}{\TMright}} (3)
    (3) edge [loop above] node {\TMtrans{\TMstroke}{\TMstroke}{\TMright}} (3)
        edge node {\TMtrans{\TMblank}{\TMstroke}{\TMright}} (4)
    (4) edge [loop below] node {\TMtrans{\TMblank}{\TMstroke}{\TMleft}} (4)
        edge node {\TMtrans{\TMstroke}{\TMstroke}{\TMleft}} (5)
    (5) edge [loop below]  node {\TMtrans{\TMstroke}{\TMstroke}{\TMleft}} (5)
        edge              node {\TMtrans{\TMblank}{\TMblank}{\TMleft}} (6)
    (6) edge [loop below] node {\TMtrans{\TMstroke}{\TMstroke}{\TMleft}} (6)
        edge              node {\TMtrans{\TMblank}{\TMblank}{\TMright}} (1);
\end{tikzpicture}
\]
\caption{آلة مضاعفة}
\ollabel{fig:doubler}
\end{figure}
\end{ex}

\begin{prob}
اختر دخلًا اعتباطيًا وتتبع تهيئات آلة المضاعفة في
\olref[tur][mac][rep]{ex:doubler}.
\end{prob}

\begin{prob}
صمم آلة تورنغ أبجديتها $\{\TMendtape,\TMblank,A,B\}$ تقبل، أي تتوقف
عند، كل سلسلة من الرمزين $A$ و$B$ يتساوى فيها عدد رموز $A$ وعدد رموز
$B$ \emph{وتسبق} فيها كل رموز $A$ جميع رموز $B$، وترفض، أي لا تتوقف
عند، كل سلسلة لا يتساوى فيها العددان أو لا تسبق فيها كل رموز $A$
جميع رموز~$B$. (ينبغي مثلًا أن تقبل الآلة $AABB$ و$AAABBB$، وأن ترفض
$AAB$ و$AABBAABB$.)
\end{prob}

\begin{prob}
صمم آلة تورنغ أبجديتها $\{\TMendtape,\TMblank,A,B\}$ تأخذ دخلًا أي
سلسلة $\alpha$ من الرمزين $A$ و$B$، وتنسخها لتنتج خرجًا على الصورة
$\alpha\alpha$. (ينبغي مثلًا أن ينتج عن الدخل $ABBA$ الخرج
$ABBAABBA$.)
\end{prob}

\begin{prob}
\emph{مرتبة أبجديًا؟} صمم آلة تورنغ أبجديتها
$\{\TMendtape,\TMblank,A,B\}$ تأخذ دخلًا متتالية منتهية من الرمزين
$A$ و$B$ وتتحقق مما إذا كانت كل رموز $A$ واقعة إلى يسار كل رموز $B$.
وينبغي أن تترك الآلة سلسلة الدخل على الشريط، وأن تتوقف إذا كانت
السلسلة «مرتبة أبجديًا»، أو تدور إلى الأبد إذا لم تكن كذلك.
\end{prob}

\begin{prob}
\emph{آلة ترتيب أبجدي:} صمم آلة تورنغ أبجديتها
$\{\TMendtape,\TMblank,A,B\}$ تأخذ دخلًا متتالية منتهية من الرمزين
$A$ و$B$ وتعيد ترتيبها بحيث تقع كل رموز $A$ إلى يسار كل رموز~$B$.
(مثلًا ينبغي أن تصير المتتالية $BABAA$ المتتالية $AAABB$، وأن تصير
المتتالية $ABBABB$ المتتالية $AABBBB$.)
\end{prob}

\end{document}
